% preamble มาตรฐานของไฟล์ข้อสอบ — ก๊อปไฟล์นี้ไปวางในโฟลเดอร์งาน แล้วบรรทัดแรกของ
% examN.tex เขียนว่า % preamble มาตรฐานของไฟล์ข้อสอบ — ก๊อปไฟล์นี้ไปวางในโฟลเดอร์งาน แล้วบรรทัดแรกของ
% examN.tex เขียนว่า % preamble มาตรฐานของไฟล์ข้อสอบ — ก๊อปไฟล์นี้ไปวางในโฟลเดอร์งาน แล้วบรรทัดแรกของ
% examN.tex เขียนว่า % preamble มาตรฐานของไฟล์ข้อสอบ — ก๊อปไฟล์นี้ไปวางในโฟลเดอร์งาน แล้วบรรทัดแรกของ
% examN.tex เขียนว่า \input{preamble-exam}
%
% ค่าทุกตัวในไฟล์นี้มีเหตุผลกำกับไว้ข้าง ๆ อย่าเปลี่ยนโดยไม่อ่านเหตุผลก่อน

\documentclass[14pt]{extarticle}
% extarticle 14pt ไม่ใช่ article 12pt — ขนาด class คุมขนาดสูตรคณิตศาสตร์ (Computer
% Modern) ถ้าลดเป็น 12pt สูตรจะเล็กกว่าตัวหนังสือไทยจนดูไม่สมดุล
% ปรับขนาดตัวหนังสือไทยที่ Scale เท่านั้น อย่าแตะ class size

\usepackage[a4paper,top=2.2cm,bottom=2.0cm,left=2.2cm,right=2.0cm,headsep=8pt]{geometry}
% ข้อสอบจริงเป็นกระดาษ Letter (612x792 pt) เพราะทำจาก Word — เราใช้ A4 โดยตั้งใจ
% เพราะปรินต์ในไทยสะดวกกว่า ถ้าผู้ใช้ขอ "เหมือนจริงทุกอย่าง" ค่อยเปลี่ยนเป็น letterpaper

\usepackage{fontspec}
\usepackage{amsmath,amssymb}
\usepackage{enumitem}
\usepackage{fancyhdr}
\usepackage{array}
\usepackage{tikz}
\usepackage{lastpage}      % ต้องคอมไพล์ 2 รอบ ไม่งั้นเลขหน้ารวมเป็น ??
\usepackage{needspace}
\usetikzlibrary{calc}

\setmainfont{THSarabunNew}[
  Extension = .ttf,
  UprightFont = *,
  BoldFont = * Bold,
  ItalicFont = * Italic,
  BoldItalicFont = * BoldItalic,
  Scale = 1.112]
% อ้างชื่อไฟล์ + Extension ไม่ใช่ชื่อ family — \setmainfont{TH Sarabun New} จะพังด้วย
% `The font "TH Sarabun New BoldItalic" cannot be found` เพราะ fontconfig ตั้งชื่อ
% หน้า BoldItalic ไม่ตรงกับที่ fontspec เดา ต้องอ้างไฟล์ตรง ๆ
%
% ไม่ต้องใส่ Path= เพราะ install.sh ก๊อปฟอนต์เข้า
% TEXMFHOME/fonts/truetype/public/thsarabunnew/ แล้ว kpathsea หาเจอเอง (portable
% ข้ามเครื่อง) ถ้าคอมไพล์แล้วหาฟอนต์ไม่เจอ ให้รัน install.sh --check ก่อน
% อย่าเติม path เต็มกลับเข้าไป
%
% Scale = 1.112 วัดมาแล้วให้ตรงข้อสอบจริงปี 68 (ฟอนต์เนื้อความ 15.95 pt) อย่าแก้ด้วยการกะ
% ถ้าต้องปรับใหม่ ดู docs/FONT.md

\XeTeXlinebreaklocale "th"
\XeTeXlinebreakskip = 0pt plus 1pt
% จำเป็นเสมอ — ไม่มีสองบรรทัดนี้ ภาษาไทยจะไม่ตัดบรรทัดเลย เพราะไม่มีช่องว่างระหว่างคำ
% XeTeX จึงมองทั้งย่อหน้าเป็นคำเดียว ข้อความจะทะลุขอบขวาออกนอกหน้า

\setlength{\parindent}{0pt}
\linespread{1.06}          % ไฟล์เฉลยทับเป็น 1.04 ใน preamble-answer.tex

\pagestyle{fancy}
\fancyhf{}
\renewcommand{\headrulewidth}{0pt}
\fancyhead[R]{\small หน้า~\thepage/\pageref{LastPage}}

% ---------------------------------------------------------------- ตัวเลือก ก/ข/ค/ง
%
% ต้องใช้ tabular + p{} ไม่ใช่ \makebox — \makebox ไม่ตัดบรรทัด ตัวเลือกยาวจะเขียนทับ
% ตัวเลือกถัดไป (เคยพลาดมาแล้วในข้อ 13 กับ 22 ชุดที่ 1) tabular ตัดบรรทัดให้แต่
% ไม่ขึ้น Overfull warning จึงต้องดูจากภาพ PNG ที่ preflight.py ชี้ให้เสมอ
%
% เลือกให้ถูก: ตัวเลือกยาวเกิน ~20 ตัวอักษรไทยใช้ \chtwo / ยาวเกินครึ่งบรรทัดใช้ \chstack

\newcommand{\chfour}[4]{%      % ตัวเลือกสั้น เรียง 4 คอลัมน์
  \par\vspace{4pt}\noindent\hspace*{1.4em}%
  \begin{tabular}{@{}p{0.20\textwidth}p{0.20\textwidth}p{0.20\textwidth}p{0.20\textwidth}@{}}
    ก.~~#1 & ข.~~#2 & ค.~~#3 & ง.~~#4
  \end{tabular}\par\vspace{3pt}}
\newcommand{\chtwo}[4]{%       % ตัวเลือกยาวปานกลาง 2x2
  \par\vspace{4pt}\noindent\hspace*{1.4em}%
  \begin{tabular}{@{}p{0.42\textwidth}p{0.42\textwidth}@{}}
    ก.~~#1 & ข.~~#2 \\[4pt]
    ค.~~#3 & ง.~~#4
  \end{tabular}\par\vspace{3pt}}
\newcommand{\chstack}[4]{%     % ตัวเลือกยาวมาก เรียงลงมา 4 บรรทัด
  \par\vspace{3pt}\noindent\hspace*{1.4em}ก.~~#1\par\vspace{1pt}%
  \noindent\hspace*{1.4em}ข.~~#2\par\vspace{1pt}%
  \noindent\hspace*{1.4em}ค.~~#3\par\vspace{1pt}%
  \noindent\hspace*{1.4em}ง.~~#4\par\vspace{2pt}}
% อาร์กิวเมนต์ทั้งสี่ของ \chstack ต้องอยู่บรรทัดเดียวกันแม้จะยาว — check_answer_key.py
% อ่านอาร์กิวเมนต์ต่อกันโดยต้องเจอ }{ ติดกัน ขึ้นบรรทัดใหม่คั่นแล้วมันจะหยุดอ่าน
% (LaTeX คอมไพล์ผ่านปกติ จึงไม่เห็นความผิดจากภาพ)

\newcommand{\note}[1]{\par\vspace{1pt}\noindent\hspace*{1.4em}\textbf{หมายเหตุ}~~#1\par}

\newcommand{\bul}[1]{\par\vspace{2pt}\noindent\hspace*{2.0em}$\bullet$~~%
  \parbox[t]{0.82\textwidth}{#1}\par\vspace{2pt}}
% ใช้ \bul แทน itemize ข้างในข้อ — check_answer_key.py นับ \item ทั้งไฟล์เป็นข้อสอบ
% ถ้ามี itemize ซ้อนอยู่ในโจทย์ เลขข้อจะเลื่อนตั้งแต่จุดนั้นเป็นต้นไป

\newlist{qlist}{enumerate}{1}
\setlist[qlist]{label=\arabic*., leftmargin=1.5em, itemsep=7pt, topsep=6pt}
% ชื่อ qlist สำคัญ — check_answer_key.py หาข้อสอบจาก qlist (ปรนัย) และ qlistb (อัตนัย)

% ใส่ \Needspace{5\baselineskip} ก่อน \item ทุกข้อ (ข้ออัตนัยใช้ 6 เฉลยใช้ 4)
% ไม่งั้นจะเจอโจทย์อยู่ท้ายหน้า แต่ตัวเลือกไปอยู่หน้าถัดไป

%
% ค่าทุกตัวในไฟล์นี้มีเหตุผลกำกับไว้ข้าง ๆ อย่าเปลี่ยนโดยไม่อ่านเหตุผลก่อน

\documentclass[14pt]{extarticle}
% extarticle 14pt ไม่ใช่ article 12pt — ขนาด class คุมขนาดสูตรคณิตศาสตร์ (Computer
% Modern) ถ้าลดเป็น 12pt สูตรจะเล็กกว่าตัวหนังสือไทยจนดูไม่สมดุล
% ปรับขนาดตัวหนังสือไทยที่ Scale เท่านั้น อย่าแตะ class size

\usepackage[a4paper,top=2.2cm,bottom=2.0cm,left=2.2cm,right=2.0cm,headsep=8pt]{geometry}
% ข้อสอบจริงเป็นกระดาษ Letter (612x792 pt) เพราะทำจาก Word — เราใช้ A4 โดยตั้งใจ
% เพราะปรินต์ในไทยสะดวกกว่า ถ้าผู้ใช้ขอ "เหมือนจริงทุกอย่าง" ค่อยเปลี่ยนเป็น letterpaper

\usepackage{fontspec}
\usepackage{amsmath,amssymb}
\usepackage{enumitem}
\usepackage{fancyhdr}
\usepackage{array}
\usepackage{tikz}
\usepackage{lastpage}      % ต้องคอมไพล์ 2 รอบ ไม่งั้นเลขหน้ารวมเป็น ??
\usepackage{needspace}
\usetikzlibrary{calc}

\setmainfont{THSarabunNew}[
  Extension = .ttf,
  UprightFont = *,
  BoldFont = * Bold,
  ItalicFont = * Italic,
  BoldItalicFont = * BoldItalic,
  Scale = 1.112]
% อ้างชื่อไฟล์ + Extension ไม่ใช่ชื่อ family — \setmainfont{TH Sarabun New} จะพังด้วย
% `The font "TH Sarabun New BoldItalic" cannot be found` เพราะ fontconfig ตั้งชื่อ
% หน้า BoldItalic ไม่ตรงกับที่ fontspec เดา ต้องอ้างไฟล์ตรง ๆ
%
% ไม่ต้องใส่ Path= เพราะ install.sh ก๊อปฟอนต์เข้า
% TEXMFHOME/fonts/truetype/public/thsarabunnew/ แล้ว kpathsea หาเจอเอง (portable
% ข้ามเครื่อง) ถ้าคอมไพล์แล้วหาฟอนต์ไม่เจอ ให้รัน install.sh --check ก่อน
% อย่าเติม path เต็มกลับเข้าไป
%
% Scale = 1.112 วัดมาแล้วให้ตรงข้อสอบจริงปี 68 (ฟอนต์เนื้อความ 15.95 pt) อย่าแก้ด้วยการกะ
% ถ้าต้องปรับใหม่ ดู docs/FONT.md

\XeTeXlinebreaklocale "th"
\XeTeXlinebreakskip = 0pt plus 1pt
% จำเป็นเสมอ — ไม่มีสองบรรทัดนี้ ภาษาไทยจะไม่ตัดบรรทัดเลย เพราะไม่มีช่องว่างระหว่างคำ
% XeTeX จึงมองทั้งย่อหน้าเป็นคำเดียว ข้อความจะทะลุขอบขวาออกนอกหน้า

\setlength{\parindent}{0pt}
\linespread{1.06}          % ไฟล์เฉลยทับเป็น 1.04 ใน preamble-answer.tex

\pagestyle{fancy}
\fancyhf{}
\renewcommand{\headrulewidth}{0pt}
\fancyhead[R]{\small หน้า~\thepage/\pageref{LastPage}}

% ---------------------------------------------------------------- ตัวเลือก ก/ข/ค/ง
%
% ต้องใช้ tabular + p{} ไม่ใช่ \makebox — \makebox ไม่ตัดบรรทัด ตัวเลือกยาวจะเขียนทับ
% ตัวเลือกถัดไป (เคยพลาดมาแล้วในข้อ 13 กับ 22 ชุดที่ 1) tabular ตัดบรรทัดให้แต่
% ไม่ขึ้น Overfull warning จึงต้องดูจากภาพ PNG ที่ preflight.py ชี้ให้เสมอ
%
% เลือกให้ถูก: ตัวเลือกยาวเกิน ~20 ตัวอักษรไทยใช้ \chtwo / ยาวเกินครึ่งบรรทัดใช้ \chstack

\newcommand{\chfour}[4]{%      % ตัวเลือกสั้น เรียง 4 คอลัมน์
  \par\vspace{4pt}\noindent\hspace*{1.4em}%
  \begin{tabular}{@{}p{0.20\textwidth}p{0.20\textwidth}p{0.20\textwidth}p{0.20\textwidth}@{}}
    ก.~~#1 & ข.~~#2 & ค.~~#3 & ง.~~#4
  \end{tabular}\par\vspace{3pt}}
\newcommand{\chtwo}[4]{%       % ตัวเลือกยาวปานกลาง 2x2
  \par\vspace{4pt}\noindent\hspace*{1.4em}%
  \begin{tabular}{@{}p{0.42\textwidth}p{0.42\textwidth}@{}}
    ก.~~#1 & ข.~~#2 \\[4pt]
    ค.~~#3 & ง.~~#4
  \end{tabular}\par\vspace{3pt}}
\newcommand{\chstack}[4]{%     % ตัวเลือกยาวมาก เรียงลงมา 4 บรรทัด
  \par\vspace{3pt}\noindent\hspace*{1.4em}ก.~~#1\par\vspace{1pt}%
  \noindent\hspace*{1.4em}ข.~~#2\par\vspace{1pt}%
  \noindent\hspace*{1.4em}ค.~~#3\par\vspace{1pt}%
  \noindent\hspace*{1.4em}ง.~~#4\par\vspace{2pt}}
% อาร์กิวเมนต์ทั้งสี่ของ \chstack ต้องอยู่บรรทัดเดียวกันแม้จะยาว — check_answer_key.py
% อ่านอาร์กิวเมนต์ต่อกันโดยต้องเจอ }{ ติดกัน ขึ้นบรรทัดใหม่คั่นแล้วมันจะหยุดอ่าน
% (LaTeX คอมไพล์ผ่านปกติ จึงไม่เห็นความผิดจากภาพ)

\newcommand{\note}[1]{\par\vspace{1pt}\noindent\hspace*{1.4em}\textbf{หมายเหตุ}~~#1\par}

\newcommand{\bul}[1]{\par\vspace{2pt}\noindent\hspace*{2.0em}$\bullet$~~%
  \parbox[t]{0.82\textwidth}{#1}\par\vspace{2pt}}
% ใช้ \bul แทน itemize ข้างในข้อ — check_answer_key.py นับ \item ทั้งไฟล์เป็นข้อสอบ
% ถ้ามี itemize ซ้อนอยู่ในโจทย์ เลขข้อจะเลื่อนตั้งแต่จุดนั้นเป็นต้นไป

\newlist{qlist}{enumerate}{1}
\setlist[qlist]{label=\arabic*., leftmargin=1.5em, itemsep=7pt, topsep=6pt}
% ชื่อ qlist สำคัญ — check_answer_key.py หาข้อสอบจาก qlist (ปรนัย) และ qlistb (อัตนัย)

% ใส่ \Needspace{5\baselineskip} ก่อน \item ทุกข้อ (ข้ออัตนัยใช้ 6 เฉลยใช้ 4)
% ไม่งั้นจะเจอโจทย์อยู่ท้ายหน้า แต่ตัวเลือกไปอยู่หน้าถัดไป

%
% ค่าทุกตัวในไฟล์นี้มีเหตุผลกำกับไว้ข้าง ๆ อย่าเปลี่ยนโดยไม่อ่านเหตุผลก่อน

\documentclass[14pt]{extarticle}
% extarticle 14pt ไม่ใช่ article 12pt — ขนาด class คุมขนาดสูตรคณิตศาสตร์ (Computer
% Modern) ถ้าลดเป็น 12pt สูตรจะเล็กกว่าตัวหนังสือไทยจนดูไม่สมดุล
% ปรับขนาดตัวหนังสือไทยที่ Scale เท่านั้น อย่าแตะ class size

\usepackage[a4paper,top=2.2cm,bottom=2.0cm,left=2.2cm,right=2.0cm,headsep=8pt]{geometry}
% ข้อสอบจริงเป็นกระดาษ Letter (612x792 pt) เพราะทำจาก Word — เราใช้ A4 โดยตั้งใจ
% เพราะปรินต์ในไทยสะดวกกว่า ถ้าผู้ใช้ขอ "เหมือนจริงทุกอย่าง" ค่อยเปลี่ยนเป็น letterpaper

\usepackage{fontspec}
\usepackage{amsmath,amssymb}
\usepackage{enumitem}
\usepackage{fancyhdr}
\usepackage{array}
\usepackage{tikz}
\usepackage{lastpage}      % ต้องคอมไพล์ 2 รอบ ไม่งั้นเลขหน้ารวมเป็น ??
\usepackage{needspace}
\usetikzlibrary{calc}

\setmainfont{THSarabunNew}[
  Extension = .ttf,
  UprightFont = *,
  BoldFont = * Bold,
  ItalicFont = * Italic,
  BoldItalicFont = * BoldItalic,
  Scale = 1.112]
% อ้างชื่อไฟล์ + Extension ไม่ใช่ชื่อ family — \setmainfont{TH Sarabun New} จะพังด้วย
% `The font "TH Sarabun New BoldItalic" cannot be found` เพราะ fontconfig ตั้งชื่อ
% หน้า BoldItalic ไม่ตรงกับที่ fontspec เดา ต้องอ้างไฟล์ตรง ๆ
%
% ไม่ต้องใส่ Path= เพราะ install.sh ก๊อปฟอนต์เข้า
% TEXMFHOME/fonts/truetype/public/thsarabunnew/ แล้ว kpathsea หาเจอเอง (portable
% ข้ามเครื่อง) ถ้าคอมไพล์แล้วหาฟอนต์ไม่เจอ ให้รัน install.sh --check ก่อน
% อย่าเติม path เต็มกลับเข้าไป
%
% Scale = 1.112 วัดมาแล้วให้ตรงข้อสอบจริงปี 68 (ฟอนต์เนื้อความ 15.95 pt) อย่าแก้ด้วยการกะ
% ถ้าต้องปรับใหม่ ดู docs/FONT.md

\XeTeXlinebreaklocale "th"
\XeTeXlinebreakskip = 0pt plus 1pt
% จำเป็นเสมอ — ไม่มีสองบรรทัดนี้ ภาษาไทยจะไม่ตัดบรรทัดเลย เพราะไม่มีช่องว่างระหว่างคำ
% XeTeX จึงมองทั้งย่อหน้าเป็นคำเดียว ข้อความจะทะลุขอบขวาออกนอกหน้า

\setlength{\parindent}{0pt}
\linespread{1.06}          % ไฟล์เฉลยทับเป็น 1.04 ใน preamble-answer.tex

\pagestyle{fancy}
\fancyhf{}
\renewcommand{\headrulewidth}{0pt}
\fancyhead[R]{\small หน้า~\thepage/\pageref{LastPage}}

% ---------------------------------------------------------------- ตัวเลือก ก/ข/ค/ง
%
% ต้องใช้ tabular + p{} ไม่ใช่ \makebox — \makebox ไม่ตัดบรรทัด ตัวเลือกยาวจะเขียนทับ
% ตัวเลือกถัดไป (เคยพลาดมาแล้วในข้อ 13 กับ 22 ชุดที่ 1) tabular ตัดบรรทัดให้แต่
% ไม่ขึ้น Overfull warning จึงต้องดูจากภาพ PNG ที่ preflight.py ชี้ให้เสมอ
%
% เลือกให้ถูก: ตัวเลือกยาวเกิน ~20 ตัวอักษรไทยใช้ \chtwo / ยาวเกินครึ่งบรรทัดใช้ \chstack

\newcommand{\chfour}[4]{%      % ตัวเลือกสั้น เรียง 4 คอลัมน์
  \par\vspace{4pt}\noindent\hspace*{1.4em}%
  \begin{tabular}{@{}p{0.20\textwidth}p{0.20\textwidth}p{0.20\textwidth}p{0.20\textwidth}@{}}
    ก.~~#1 & ข.~~#2 & ค.~~#3 & ง.~~#4
  \end{tabular}\par\vspace{3pt}}
\newcommand{\chtwo}[4]{%       % ตัวเลือกยาวปานกลาง 2x2
  \par\vspace{4pt}\noindent\hspace*{1.4em}%
  \begin{tabular}{@{}p{0.42\textwidth}p{0.42\textwidth}@{}}
    ก.~~#1 & ข.~~#2 \\[4pt]
    ค.~~#3 & ง.~~#4
  \end{tabular}\par\vspace{3pt}}
\newcommand{\chstack}[4]{%     % ตัวเลือกยาวมาก เรียงลงมา 4 บรรทัด
  \par\vspace{3pt}\noindent\hspace*{1.4em}ก.~~#1\par\vspace{1pt}%
  \noindent\hspace*{1.4em}ข.~~#2\par\vspace{1pt}%
  \noindent\hspace*{1.4em}ค.~~#3\par\vspace{1pt}%
  \noindent\hspace*{1.4em}ง.~~#4\par\vspace{2pt}}
% อาร์กิวเมนต์ทั้งสี่ของ \chstack ต้องอยู่บรรทัดเดียวกันแม้จะยาว — check_answer_key.py
% อ่านอาร์กิวเมนต์ต่อกันโดยต้องเจอ }{ ติดกัน ขึ้นบรรทัดใหม่คั่นแล้วมันจะหยุดอ่าน
% (LaTeX คอมไพล์ผ่านปกติ จึงไม่เห็นความผิดจากภาพ)

\newcommand{\note}[1]{\par\vspace{1pt}\noindent\hspace*{1.4em}\textbf{หมายเหตุ}~~#1\par}

\newcommand{\bul}[1]{\par\vspace{2pt}\noindent\hspace*{2.0em}$\bullet$~~%
  \parbox[t]{0.82\textwidth}{#1}\par\vspace{2pt}}
% ใช้ \bul แทน itemize ข้างในข้อ — check_answer_key.py นับ \item ทั้งไฟล์เป็นข้อสอบ
% ถ้ามี itemize ซ้อนอยู่ในโจทย์ เลขข้อจะเลื่อนตั้งแต่จุดนั้นเป็นต้นไป

\newlist{qlist}{enumerate}{1}
\setlist[qlist]{label=\arabic*., leftmargin=1.5em, itemsep=7pt, topsep=6pt}
% ชื่อ qlist สำคัญ — check_answer_key.py หาข้อสอบจาก qlist (ปรนัย) และ qlistb (อัตนัย)

% ใส่ \Needspace{5\baselineskip} ก่อน \item ทุกข้อ (ข้ออัตนัยใช้ 6 เฉลยใช้ 4)
% ไม่งั้นจะเจอโจทย์อยู่ท้ายหน้า แต่ตัวเลือกไปอยู่หน้าถัดไป

%
% ค่าทุกตัวในไฟล์นี้มีเหตุผลกำกับไว้ข้าง ๆ อย่าเปลี่ยนโดยไม่อ่านเหตุผลก่อน

\documentclass[14pt]{extarticle}
% extarticle 14pt ไม่ใช่ article 12pt — ขนาด class คุมขนาดสูตรคณิตศาสตร์ (Computer
% Modern) ถ้าลดเป็น 12pt สูตรจะเล็กกว่าตัวหนังสือไทยจนดูไม่สมดุล
% ปรับขนาดตัวหนังสือไทยที่ Scale เท่านั้น อย่าแตะ class size

\usepackage[a4paper,top=2.2cm,bottom=2.0cm,left=2.2cm,right=2.0cm,headsep=8pt]{geometry}
% ข้อสอบจริงเป็นกระดาษ Letter (612x792 pt) เพราะทำจาก Word — เราใช้ A4 โดยตั้งใจ
% เพราะปรินต์ในไทยสะดวกกว่า ถ้าผู้ใช้ขอ "เหมือนจริงทุกอย่าง" ค่อยเปลี่ยนเป็น letterpaper

\usepackage{fontspec}
\usepackage{amsmath,amssymb}
\usepackage{enumitem}
\usepackage{fancyhdr}
\usepackage{array}
\usepackage{tikz}
\usepackage{lastpage}      % ต้องคอมไพล์ 2 รอบ ไม่งั้นเลขหน้ารวมเป็น ??
\usepackage{needspace}
\usetikzlibrary{calc}

\setmainfont{THSarabunNew}[
  Extension = .ttf,
  UprightFont = *,
  BoldFont = * Bold,
  ItalicFont = * Italic,
  BoldItalicFont = * BoldItalic,
  Scale = 1.112]
% อ้างชื่อไฟล์ + Extension ไม่ใช่ชื่อ family — \setmainfont{TH Sarabun New} จะพังด้วย
% `The font "TH Sarabun New BoldItalic" cannot be found` เพราะ fontconfig ตั้งชื่อ
% หน้า BoldItalic ไม่ตรงกับที่ fontspec เดา ต้องอ้างไฟล์ตรง ๆ
%
% ไม่ต้องใส่ Path= เพราะ install.sh ก๊อปฟอนต์เข้า
% TEXMFHOME/fonts/truetype/public/thsarabunnew/ แล้ว kpathsea หาเจอเอง (portable
% ข้ามเครื่อง) ถ้าคอมไพล์แล้วหาฟอนต์ไม่เจอ ให้รัน install.sh --check ก่อน
% อย่าเติม path เต็มกลับเข้าไป
%
% Scale = 1.112 วัดมาแล้วให้ตรงข้อสอบจริงปี 68 (ฟอนต์เนื้อความ 15.95 pt) อย่าแก้ด้วยการกะ
% ถ้าต้องปรับใหม่ ดู docs/FONT.md

\XeTeXlinebreaklocale "th"
\XeTeXlinebreakskip = 0pt plus 1pt
% จำเป็นเสมอ — ไม่มีสองบรรทัดนี้ ภาษาไทยจะไม่ตัดบรรทัดเลย เพราะไม่มีช่องว่างระหว่างคำ
% XeTeX จึงมองทั้งย่อหน้าเป็นคำเดียว ข้อความจะทะลุขอบขวาออกนอกหน้า

\setlength{\parindent}{0pt}
\linespread{1.06}          % ไฟล์เฉลยทับเป็น 1.04 ใน preamble-answer.tex

\pagestyle{fancy}
\fancyhf{}
\renewcommand{\headrulewidth}{0pt}
\fancyhead[R]{\small หน้า~\thepage/\pageref{LastPage}}

% ---------------------------------------------------------------- ตัวเลือก ก/ข/ค/ง
%
% ต้องใช้ tabular + p{} ไม่ใช่ \makebox — \makebox ไม่ตัดบรรทัด ตัวเลือกยาวจะเขียนทับ
% ตัวเลือกถัดไป (เคยพลาดมาแล้วในข้อ 13 กับ 22 ชุดที่ 1) tabular ตัดบรรทัดให้แต่
% ไม่ขึ้น Overfull warning จึงต้องดูจากภาพ PNG ที่ preflight.py ชี้ให้เสมอ
%
% เลือกให้ถูก: ตัวเลือกยาวเกิน ~20 ตัวอักษรไทยใช้ \chtwo / ยาวเกินครึ่งบรรทัดใช้ \chstack

\newcommand{\chfour}[4]{%      % ตัวเลือกสั้น เรียง 4 คอลัมน์
  \par\vspace{4pt}\noindent\hspace*{1.4em}%
  \begin{tabular}{@{}p{0.20\textwidth}p{0.20\textwidth}p{0.20\textwidth}p{0.20\textwidth}@{}}
    ก.~~#1 & ข.~~#2 & ค.~~#3 & ง.~~#4
  \end{tabular}\par\vspace{3pt}}
\newcommand{\chtwo}[4]{%       % ตัวเลือกยาวปานกลาง 2x2
  \par\vspace{4pt}\noindent\hspace*{1.4em}%
  \begin{tabular}{@{}p{0.42\textwidth}p{0.42\textwidth}@{}}
    ก.~~#1 & ข.~~#2 \\[4pt]
    ค.~~#3 & ง.~~#4
  \end{tabular}\par\vspace{3pt}}
\newcommand{\chstack}[4]{%     % ตัวเลือกยาวมาก เรียงลงมา 4 บรรทัด
  \par\vspace{3pt}\noindent\hspace*{1.4em}ก.~~#1\par\vspace{1pt}%
  \noindent\hspace*{1.4em}ข.~~#2\par\vspace{1pt}%
  \noindent\hspace*{1.4em}ค.~~#3\par\vspace{1pt}%
  \noindent\hspace*{1.4em}ง.~~#4\par\vspace{2pt}}
% อาร์กิวเมนต์ทั้งสี่ของ \chstack ต้องอยู่บรรทัดเดียวกันแม้จะยาว — check_answer_key.py
% อ่านอาร์กิวเมนต์ต่อกันโดยต้องเจอ }{ ติดกัน ขึ้นบรรทัดใหม่คั่นแล้วมันจะหยุดอ่าน
% (LaTeX คอมไพล์ผ่านปกติ จึงไม่เห็นความผิดจากภาพ)

\newcommand{\note}[1]{\par\vspace{1pt}\noindent\hspace*{1.4em}\textbf{หมายเหตุ}~~#1\par}

\newcommand{\bul}[1]{\par\vspace{2pt}\noindent\hspace*{2.0em}$\bullet$~~%
  \parbox[t]{0.82\textwidth}{#1}\par\vspace{2pt}}
% ใช้ \bul แทน itemize ข้างในข้อ — check_answer_key.py นับ \item ทั้งไฟล์เป็นข้อสอบ
% ถ้ามี itemize ซ้อนอยู่ในโจทย์ เลขข้อจะเลื่อนตั้งแต่จุดนั้นเป็นต้นไป

\newlist{qlist}{enumerate}{1}
\setlist[qlist]{label=\arabic*., leftmargin=1.5em, itemsep=7pt, topsep=6pt}
% ชื่อ qlist สำคัญ — check_answer_key.py หาข้อสอบจาก qlist (ปรนัย) และ qlistb (อัตนัย)

% ใส่ \Needspace{5\baselineskip} ก่อน \item ทุกข้อ (ข้ออัตนัยใช้ 6 เฉลยใช้ 4)
% ไม่งั้นจะเจอโจทย์อยู่ท้ายหน้า แต่ตัวเลือกไปอยู่หน้าถัดไป
